\chapter{Référentiels et Méthodes d'Analyse}
\label{chap:referentiels}

L'analyse de risque ne s'improvise pas. Pour être reconnue et cohérente, elle doit suivre des cadres méthodologiques précis, appelés \textbf{référentiels}.

\section{Le Paysage Normatif}

\subsection{Qu'est-ce qu'une norme ISO ?}

\begin{boxdef}
\textbf{ISO} : Acronyme de \textit{International Organization for Standardization}. C'est un organisme international qui définit des standards volontaires. Une norme ISO est un document qui fournit des exigences, des spécifications, des lignes directrices ou des caractéristiques à utiliser de manière cohérente.
\end{boxdef}

\subsection{ISO 31000 : La Norme "Mère"}

La norme \textbf{ISO 31000} est le standard international de référence pour le management du risque. Elle ne se limite pas à l'informatique ; elle s'applique à tout type de risque (financier, sécurité physique, projet).

\textbf{Principes clés :}
\begin{itemize}
    \item Le management du risque doit être intégré, structuré et inclusif.
    \item Il doit être dynamique (le risque évolue) et adapté au contexte.
\end{itemize}

Elle fournit le cadre général, mais pas les outils techniques "prêts à l'emploi" pour le risque numérique. C'est là qu'intervient l'ISO 27005.

\subsection{ISO 27005 : Spécialisée Sécurité de l'Information}

La norme \textbf{ISO 27005} (Sécurité de l'information -- Gestion des risques en sécurité de l'information) est le dérivé de l'ISO 31000 pour le monde du SI.

\begin{boxdef}
\textbf{Objectif ISO 27005 :} Fournir des lignes directrices pour la gestion des risques liés à la sécurité de l'information. Elle soutient les exigences de l'ISO 27001.
\end{boxdef}

Elle définit un processus itératif :
\begin{enumerate}
    \item \textbf{Contextualisation} : Définir le cadre.
    \item \textbf{Appréciation du risque} : Identification, Analyse, Estimation.
    \item \textbf{Traitement du risque} : Réduction, Évitement, Transfert, Acceptation.
    \item \textbf{Acceptation du risque} : Validation formelle.
    \item \textbf{Communication} : Dialogue avec les parties prenantes.
    \item \textbf{Suivi et révision} : Mise à jour continue.
\end{enumerate}

% ------------------------------------------------------------------------------
\section{EBIOS RM : La Méthode de Référence Française}
% ------------------------------------------------------------------------------

\subsection{Présentation et Origine}

\begin{boxdef}
\textbf{EBIOS RM} : \textit{Expression des Besoins et Identification des Objectifs de Sécurité - Risk Manager}.
C'est une méthode d'analyse de risques développée par l'ANSSI (\textit{Agence Nationale de la Sécurité des Systèmes d'Information}) en France.
\end{boxdef}

C'est la méthode la plus utilisée en France par les administrations et les grandes entreprises privées, notamment pour sa conformité avec le RGPD et la Loi de Programmation Militaire (LPM).

\subsection{Les 5 Ateliers de la Méthode}

EBIOS RM fonctionne par ateliers collaboratifs. Les participants (métiers et IT) travaillent ensemble.

\subsubsection{Atelier 1 : Établissement du contexte}
On définit le périmètre, les contraintes (budgétaires, légales) et on identifie les \textbf{Actifs Essentiels} (ce qui a de la valeur).
\begin{boxlizbiz}
\textbf{LiZbiZ :} L'actif essentiel identifié est la "Base de Données Clients" et le "Processus de Commande".
\end{boxlizbiz}

\subsubsection{Atelier 2 : Étude des évènements redoutés}
On identifie les \textbf{Menaces} et les \textbf{Sources de Menaces}. On classe les risques selon les critères de sécurité :
\begin{itemize}
    \item \textbf{D} : Disponibilité.
    \item \textbf{I} : Intégrité.
    \item \textbf{C} : Confidentialité.
    \item \textbf{T} : Traçabilité (Preuve).
\end{itemize}

\subsubsection{Atelier 3 : Analyse des risques}
On identifie les \textbf{Vulnérabilités} et on construit des \textbf{Scénarios de Risque}.
Un scénario répond à la question : \textit{"Une source de menace exploite-t-elle une vulnérabilité pour nuire à un actif ?"}

\subsubsection{Atelier 4 : Étude des mesures}
On évalue la gravité et la vraisemblance pour calculer le risque. On identifie les mesures existantes et leur efficacité.

\subsubsection{Atelier 5 : Gestion du risque}
On décide de la stratégie de traitement (Accepter, Réduire, etc.) et on élabore le \textbf{Plan d'Action Sécurité (PAS)}.

% ------------------------------------------------------------------------------
\section{NIST SP 800-30 : L'Approche Américaine}
% ------------------------------------------------------------------------------

\subsection{Présentation}

Le \textbf{NIST} (\textit{National Institute of Standards and Technology}) est un organisme américain. Le document \textbf{SP 800-30} (Guide for Conducting Risk Assessments) est le standard de fait aux États-Unis et pour les entreprises travaillant avec le gouvernement américain.

\subsection{Comparaison avec EBIOS}

\begin{table}[htbp]
\centering
\begin{tabular}{p{3cm}p{5cm}p{5cm}}
    \toprule
    \textbf{Critère} & \textbf{EBIOS RM (France)} & \textbf{NIST SP 800-30 (USA)} \\
    \midrule
    Approche & Collaborative, ateliers humains. & Processus,Base de données technique. \\
    Focus & Gouvernance et métier. & Technique et opérationnel. \\
    Format & Grille de probabilité/gravité (Qualitatif). & Grille High/Medium/Low (Qualitatif/Semi-quantitatif). \\
    Terminologie & "Evènements redoutés". & "Threat Events". \\
    \bottomrule
\end{tabular}
\caption{Comparaison EBIOS RM vs NIST SP 800-30}
\end{table}

\begin{boxdef}
\textbf{RMF (Risk Management Framework) :} Le cadre du NIST qui inclut la gestion du risque dans tout le cycle de vie du système (Categorize, Select, Implement, Assess, Authorize, Monitor).
\end{boxdef}

% ------------------------------------------------------------------------------
\section{Cas LiZbiZ : Choix de la Méthode}
% ------------------------------------------------------------------------------

\subsection{Analyse du besoin méthodologique}

LiZbiZ est une entreprise française en forte croissance, soumise au RGPD et potentiellement à la LPM (secteur logistique). Le PDG est soucieux de sa réputation et la DSI a besoin d'une méthode structurée pour communiquer avec le COMEX.

\subsection{Décision}

\begin{boxlizbiz}
\textbf{Choix retenu : La méthode EBIOS RM.}

\textbf{Justification :}
\begin{enumerate}
    \item \textbf{Conformité :} Alignement total avec les exigences réglementaires françaises et européennes.
    \item \textbf{Communication :} Les scénarios d'EBIOS sont compréhensibles par les non-techniciens (COMEX), ce qui est vital pour obtenir des budgets.
    \item \textbf{Gouvernance :} Intègre naturellement les concepts de RACI et de validation par la direction.
\end{enumerate}

Nous utiliserons EBIOS RM pour l'ensemble des ateliers pratiques de ce module.
\end{boxlizbiz}


% ==============================================================================
% Fichier : 02_module_risque.tex (Extrait Chapitre 4)
% Description : Application Pratique - Analyse de Risque LiZbiZ
% ==============================================================================